\documentclass{situationist-preprint}

\title{Strawson's Objection Reconsidered}
\msid{LNG-03-b}
\msfamily{LNG-03}
\msdiscipline{Philosophy of Language}
\mstype{Consequence}
\mskeywords{definite descriptions; presupposition; Russell; Strawson; truth-value gaps; pragmatics}

\begin{document}
\maketitle

\begin{abstract}
LNG-03-a proves that, on Russell's analysis, \enquote{The present King of France is bald} is false, and claims that Strawson's objection concerns use rather than content. We derive the second claim from the first together with a distinction between the proposition an utterance expresses and the conditions under which uttering it is in order. Each of Strawson's observations, that a hearer would not answer \enquote{false}, that the question of truth seems not to arise, and that the speaker has failed to refer, is a datum about the use of the sentence, and each is predicted by the failure of a pragmatic presupposition. None requires a truth-value gap. The asymmetries Strawson himself later noted confirm the division.
\end{abstract}

\section{The objection}

Strawson argued that a speaker who uttered \enquote{The King of France is wise} when there was no king of France would not have said anything false; the question whether the statement was true or false would simply not arise, because the speaker would have failed to mention anyone \citep{Strawson1950}. His positive account distinguished the sentence, which is meaningful whether or not there is a king, from its use on an occasion to make a statement, and held that the use presupposes, without asserting, that there is a king.

LNG-03-a proves (Theorem~3.2) that on Russell's analysis \citep{Russell1905}, given that nothing is presently king of France (Assumption~2.2), the sentence is false, and (Corollary~3.3) that its narrow-scope negation is false and its wide-scope negation true. The question is how to square these verdicts with Strawson's observations. Our answer is that the observations are correct and concern a different object.

\section{Content and use}

\begin{definition}\label{d:content}
The \emph{content} of an utterance is the proposition it expresses, which is true or false according to how things are.
\end{definition}

\begin{definition}\label{d:presup}
A speaker \emph{presupposes} that $p$ in making an utterance if she takes $p$ for granted and acts as if she takes her audience to take it for granted \citep{Stalnaker1974}. An utterance is \emph{in order} only if its speaker's presuppositions are true.
\end{definition}

The two notions are independent. Content is a property of what is said; presupposition is a property of the speaker's conduct in saying it. An utterance can have a false content and be in order (an honest mistake), and it can have a false content and be out of order because a presupposition fails.

\begin{proposition}\label{p:use}
An utterance of \enquote{The present King of France is bald}, made by a speaker who takes it for granted that there is a king of France, has a false content and is not in order.
\end{proposition}

\begin{proof}
The content is false by Theorem~3.2 of LNG-03-a. The speaker presupposes that there is a king of France, and the presupposition is false by Assumption~2.2 of LNG-03-a; so the utterance is not in order by Definition~\ref{d:presup}.
\end{proof}

Strawson's observations are now explained. A hearer asked to respond does not answer \enquote{false} because a bare \enquote{false} would be taken as accepting the shared presupposition and denying only the baldness, that is, as asserting the narrow-scope negation, which is itself false by Corollary~3.3 of LNG-03-a. The cooperative response targets the presupposition: \enquote{There is no king of France.} The impression that the question of truth does not arise is the impression that the utterance is out of order, and so is not one to be assessed on its merits. The claim that the speaker has failed to mention anyone is true and concerns the speaker's act of referring, which belongs to use. None of the three observations requires that the content lack a truth value.

\begin{corollary}\label{c:division}
Strawson's data are facts about the presuppositions of the use of the sentence; Russell's verdicts are facts about its content. Both can be correct.
\end{corollary}

\section{Strawson's later asymmetries}

Strawson later observed that judgements vary with topic \citep{Strawson1964}. Asked about \enquote{The King of France is bald}, hearers are reluctant to say \enquote{false}; told that the exhibition was visited yesterday by the King of France, they say without hesitation that the statement is false. Von Fintel found the same pattern and proposed that hearers can reject an utterance with a failed presupposition when it makes a claim that can be refuted independently of the failure \citep{vonFintel2004}.

On the present division the asymmetry is expected. In the exhibition case, the description is not the topic; the speaker presupposes only that the exhibition took place, which is true, and the utterance is in order with respect to its topic. Nothing then interferes with the assessment of the content, and the content is judged false. In the baldness case the description is the topic and the presupposition about it fails, so the response targets the presupposition. If content had a gap whenever a description is empty, both utterances would lack a truth value, and the confident verdict in the exhibition case would be an error. The theory that locates Strawson's observations in use predicts both judgements; the gap theory predicts only one.

\begin{objection}
If the presupposition is not part of the content, a speaker who utters the sentence asserts nothing about there being a king. Yet the hearer can protest \enquote{You said that there is a king of France}.
\end{objection}

\begin{reply}
On Russell's analysis the content entails that there is a king of France, so the speaker did say it, as part of a single existential proposition. What is presupposed is that the audience already accepts the existential conjunct. The protest is apt on both counts: the hearer rejects a conjunct of what was said and the presumption that it was common ground.
\end{reply}

\begin{objection}
The reply makes the use/content distinction carry all the weight. A gap theorist could equally say that the exhibition judgement concerns use.
\end{objection}

\begin{reply}
The gap theorist would have to say that hearers judge \enquote{false} an utterance whose content is neither true nor false, and so that a confident, stable and widely shared judgement of falsity is mistaken. The Russellian takes the judgement at face value and explains the reluctance in the other case by a recognized mechanism. Other things equal, the account that does not convict competent speakers of systematic error is preferable.
\end{reply}

\section{Conclusion}

From the falsity of the sentence, established in LNG-03-a, and the independence of content and presupposition, it follows that Strawson's objection concerns the use of the sentence. What Strawson saw, that the utterance is defective and that a plain \enquote{false} is the wrong response, is true; it is a fact about what the speaker takes for granted, and it leaves the content false.

\begin{thebibliography}{9}\small
\bibitem[Russell(1905)]{Russell1905} Russell, B. (1905). On denoting. \emph{Mind}, 14, 479--493.
\bibitem[Stalnaker(1974)]{Stalnaker1974} Stalnaker, R. (1974). Pragmatic presuppositions. In M.~K. Munitz and P.~K. Unger (eds.), \emph{Semantics and Philosophy}. New York University Press.
\bibitem[Strawson(1950)]{Strawson1950} Strawson, P.~F. (1950). On referring. \emph{Mind}, 59, 320--344.
\bibitem[Strawson(1964)]{Strawson1964} Strawson, P.~F. (1964). Identifying reference and truth-values. \emph{Theoria}, 30, 96--118.
\bibitem[von Fintel(2004)]{vonFintel2004} von Fintel, K. (2004). Would you believe it? The King of France is back! (Presuppositions and truth-value intuitions). In M.~Reimer and A.~Bezuidenhout (eds.), \emph{Descriptions and Beyond}. Oxford University Press.
\end{thebibliography}

\end{document}
